\scanpage{26}
% archive sheet number: 25
% Source lines 1--2.
\uline{Théorème} $G$ désigne un groupe affine connexe sur $k$, $B$ Borel, $\unsure{\mathscr L|\mathcal L}\subset\unsure{\mathfrak g|\mathscr G}$ alg. de Lie. Alors
\[
B=\operatorname{Norm}_{G}(B)\tag*{a)}
\]
\margin{? ---}
\[
\begin{aligned}
B&\longrightarrow\operatorname{Norm}_{G}(\mathscr L)\quad\text{est \uline{surjectif}, donc}\\
\operatorname{Drap}(G)&\longrightarrow\operatorname{Grass}(\mathfrak g)\quad\text{injectif (pas tjrs une imm.)}
\end{aligned}
\tag*{b)}
\]

% Source lines 6--9; an earlier item letter is struck above c).
\deleted{\unsure{c)|b)}} c) Conditions équivalentes
\[
\begin{aligned}
\text{i)}\quad B&=\operatorname{Norm}_{G}(\mathscr L)
&\text{ii)}\quad\mathscr L&=\operatorname{Norm}_{\mathfrak g}(\mathscr L)\\
\text{iii)}\quad\operatorname{Drap}(G)&\longrightarrow\operatorname{Grass}(\mathfrak g)\quad\text{une imm.}
\end{aligned}
\]
\smallskip\hrule\smallskip
% Source lines 10--13.
On peut supposer $k$ alg. clos

 a) Réduction au cas $\unsure{k|R}$ \unsure{simple} ---

 b) \deleted{\unsure{cas}} Réduction au cas $\unsure{k|R}$ \unsure{simple} \deleted{\unsure{un}} \unsure{un.}

\deleted{\unsure{Soit} $U=B_{\unsure{u|u.a}}$}
